\documentclass[main.tex]{subfiles}


\begin{document}

% \AddToShipoutPictureFG{%
% \begin{tikzpicture}[overlay,remember picture]
% \path (current page.south west) -- (current page.north east)
% node[midway,scale=1,color=gray,sloped,opacity=0.4] {\textnormal{Кравченко Игорь Игоревич\hspace{1cm} +79010144910\hspace{1cm} super.antig@yandex.ru}};
% \end{tikzpicture}
% }


% %from pagenumber to up
% \phantomsection 
% \hypertarget{toc}{}

 

 

% номер секции вручную
\setcounter{section}{22
}
\addtocounter{section}{-1} 

\section{Давление}

В ряде случаев существенно то, что силы, действующие на тела, распределены по некоторой части поверхности этих тел. На рис.\,\ref{fig:p42} изображены два одинаковых бруска, покоящихся на горизонтальной поверхности.
% \footnote{Интерактивные изображения ниже можно вращать мышью в программе Adobe Reader.}.


\def\sub{.417}
\begin{figure}[H]%
\hfill
\begin{subfigure}{\sub\linewidth}
\centering
\includegraphics[]{mech/30press_pic1.pdf}%
% \caption{При $s_1=s_2$ вторая машина затратила меньше времени.}
% \label{fig:p42sub1}
\end{subfigure}%
\hfill
\begin{subfigure}{\sub\linewidth}
\centering
\includegraphics[]{mech/30press_pic2.pdf}%
% \caption{При $t_1=t_2$ обе машины прошли одинаковое расстояние.}
% \label{fig:p42sub2}
\end{subfigure}%
\hfill\mbox{}
\caption{Бруски на поверхности}
\label{fig:p42}
\end{figure}


Бруски~Б$_1$ и~Б$_2$ соприкасаются с опорой поверхностями площадями~$S_1$ и~$S_2$ соответственно ($S_2>S_1$). При этом, так как массы брусков одинаковы, возникают одинаковые приложенные к опоре силы (вес), обозначаемые~$\vec F$ для общего случая (рис.\,\ref{fig:p43}). 

%%%%%%%%%%%%%%%%%%%%%%%%%%%%%%%%%%%%%
%next pic

\def\sub{.417}
\begin{figure}[H]%
\hfill
\begin{subfigure}{\sub\linewidth}
\centering
\includegraphics[]{mech/30press_pic3.pdf}%
% \caption{При $s_1=s_2$ вторая машина затратила меньше времени.}
% \label{fig:p42sub1}
\end{subfigure}%
\hfill
\begin{subfigure}{\sub\linewidth}
\centering
\includegraphics[]{mech/30press_pic4.pdf}%
% \caption{При $t_1=t_2$ обе машины прошли одинаковое расстояние.}
% \label{fig:p42sub2}
\end{subfigure}%
\hfill\mbox{}
\caption{Силы и соответствующие площади}
\label{fig:p43}
\end{figure}

Если поверхности опор достаточно мягкие, то на практике тела оставляют довольно заметные вмятины под собой. При тонких измерения <<погружения>> в опору брусков на рис.\,\ref{fig:p42} оказывается, что глубина, на которую <<вошел>> брусок~Б$_1$, больше соответствующей глубины для бруска~Б$_2$ (несмотря на то, что силы, приложенные к опоре со стороны каждого из брусков, одинаковы, что проиллюстрировано на рис.\,\ref{fig:p43}!).

\textbf{Давление} ($P_\text{д}$ [Па])\footnote{Подпись~<<д>> необязательна и введена для того, чтобы не путать с весом, обозначаемым~$\vec P$.}~--- это характеристика <<продавливающего>> действия силы:
\begin{equation}\label{30press_e1}
    P_\text{д}=\frac{F_\perp}{S},
\end{equation}
где~$F_\perp$~--- перпендикулярная составляющая силы, действующей на поверхность площади~$S$.

Также можно сказать, что давление~--- это \emph{поверхностная плотность силы}, то есть эта величина показывает, какая сила приходится на единицу площади.

Возвращаясь теперь к рис.\,\ref{fig:p43}, можно видеть, что согласно формуле~\eqref{30press_e1} давления~$P_1$ и~$P_2$, создаваемые брусками~Б$_1$ и~Б$_2$ соответственно, связаны следующим соотношением: $P_1>P_2$.











 
\end{document}